\part{Automatismes et évaluations}
\chapter{Automatismes de chaque séance}

Chaque série est prévue pour 8 à 10 minutes, sans calculatrice sauf indication.
Elle associe la notion du jour à une question de rappel. Dans l'édition
professeur, les réponses apparaissent immédiatement sous la série.

\section{Semaine 1, séance A -- Énoncés, propositions et vérité}
\begin{enumerate}
  \item Dire si « $7$ est pair » est une proposition et donner sa valeur de vérité.
  \item Nier : « $x\geq 3$ ».
  \item Nier : « $P$ et $Q$ ».
  \item Nier : « $P$ ou $Q$ ».
  \item Encadrer $\sqrt{20}$ entre deux entiers consécutifs.
\end{enumerate}
\TeacherOnly{%
\begin{tcolorbox}[colback=VE!3,colframe=VE!35,boxrule=.4pt,arc=1mm]
\small\textit{Réponses.}\quad \textbf{1.} Proposition fausse. \quad \textbf{2.} $x<3$. \quad \textbf{3.} « non $P$ ou non $Q$ ». \quad \textbf{4.} « non $P$ et non $Q$ ». \quad \textbf{5.} $4<\sqrt{20}<5$.
\end{tcolorbox}}

\section{Semaine 1, séance B -- Connecteurs et négation}
\begin{enumerate}
  \item Donner la réciproque de $P\Rightarrow Q$.
  \item Donner la contraposée de $P\Rightarrow Q$.
  \item Traduire « tout réel a un carré positif ou nul ».
  \item Nier : $\forall x\in\R,\ x^2\geq1$.
  \item Calculer $\dfrac34-\dfrac16$.
\end{enumerate}
\TeacherOnly{%
\begin{tcolorbox}[colback=VE!3,colframe=VE!35,boxrule=.4pt,arc=1mm]
\small\textit{Réponses.}\quad \textbf{1.} $Q\Rightarrow P$. \quad \textbf{2.} $\neg Q\Rightarrow\neg P$. \quad \textbf{3.} $\forall x\in\R,\ x^2\geq0$. \quad \textbf{4.} $\exists x\in\R,\ x^2<1$. \quad \textbf{5.} $\dfrac7{12}$.
\end{tcolorbox}}

\section{Semaine 2, séance A -- Implication et équivalence}
\begin{enumerate}
  \item La phrase « $x+2=5$ » est-elle une proposition lorsque $x$ n'est pas fixé ?
  \item Compléter : $P\Leftrightarrow Q$ signifie $P\Rightarrow Q$ et \ldots
  \item Donner un contre-exemple à « tout nombre premier est impair ».
  \item Si $P$ est vraie et $Q$ fausse, donner la valeur de $P\land Q$.
  \item Calculer $3^2\times3^{-4}$.
\end{enumerate}
\TeacherOnly{%
\begin{tcolorbox}[colback=VE!3,colframe=VE!35,boxrule=.4pt,arc=1mm]
\small\textit{Réponses.}\quad \textbf{1.} Non : c'est un prédicat. \quad \textbf{2.} $Q\Rightarrow P$. \quad \textbf{3.} $2$. \quad \textbf{4.} Fausse. \quad \textbf{5.} $\dfrac19$.
\end{tcolorbox}}

\section{Semaine 2, séance B -- Quantificateurs et négations}
\begin{enumerate}
  \item Si $P$ est vraie et $Q$ fausse, donner la valeur de $P\lor Q$.
  \item Nier : « il existe un entier pair supérieur à $10$ ».
  \item Écrire « $2$ est l'unique solution » avec le symbole $\exists!$.
  \item La négation de $x\neq0$ est-elle $x=0$ ?
  \item Donner $]-\infty;3]\cup[1;6[$.
\end{enumerate}
\TeacherOnly{%
\begin{tcolorbox}[colback=VE!3,colframe=VE!35,boxrule=.4pt,arc=1mm]
\small\textit{Réponses.}\quad \textbf{1.} Vraie. \quad \textbf{2.} Tout entier supérieur à $10$ est impair. \quad \textbf{3.} $\exists!x,\ x=2$. \quad \textbf{4.} Oui. \quad \textbf{5.} $]-\infty;6[$.
\end{tcolorbox}}

\section{Semaine 3, séance A -- Variables et affectation}
\begin{enumerate}
  \item Après $x\leftarrow4$ puis $x\leftarrow2x-1$, que vaut $x$ ?
  \item Après $a\leftarrow3$, $b\leftarrow5$, $a\leftarrow b$, que vaut $a$ ?
  \item Que teste la condition $n\bmod2=0$ ?
  \item Combien de fois s'exécute une boucle « Pour $i$ de $1$ à $6$ » ?
  \item Comparer $\sqrt{11}$ et $\dfrac72$.
\end{enumerate}
\TeacherOnly{%
\begin{tcolorbox}[colback=VE!3,colframe=VE!35,boxrule=.4pt,arc=1mm]
\small\textit{Réponses.}\quad \textbf{1.} $7$. \quad \textbf{2.} $5$. \quad \textbf{3.} Si $n$ est pair. \quad \textbf{4.} $6$ fois. \quad \textbf{5.} $\sqrt{11}<\dfrac72$.
\end{tcolorbox}}

\section{Semaine 3, séance B -- Lire un pseudo-code}
\begin{enumerate}
  \item Avec $S\leftarrow0$, puis $S\leftarrow S+3$ répété quatre fois, donner $S$.
  \item Dans un organigramme, quelle forme représente un test ?
  \item Dans un organigramme, quelle forme représente un traitement ?
  \item Après $x\leftarrow-2$, afficher $x^2+1$.
  \item Résoudre $|x|=6$.
\end{enumerate}
\TeacherOnly{%
\begin{tcolorbox}[colback=VE!3,colframe=VE!35,boxrule=.4pt,arc=1mm]
\small\textit{Réponses.}\quad \textbf{1.} $12$. \quad \textbf{2.} Un losange. \quad \textbf{3.} Un rectangle. \quad \textbf{4.} $5$. \quad \textbf{5.} $x=-6$ ou $x=6$.
\end{tcolorbox}}

\section{Semaine 4, séance A -- Conditions}
\begin{enumerate}
  \item Si $x=7$, le test $x<5$ est-il vrai ?
  \item Avec $i\leftarrow1$, puis tant que $i<5$, $i\leftarrow i+1$, donner la valeur finale.
  \item Différencier en un mot « Lire $x$ » et « Afficher $x$ ».
  \item Après $a\leftarrow2$, $b\leftarrow a+3$, $a\leftarrow b^2$, donner $(a,b)$.
  \item Calculer $(2^3)^2$.
\end{enumerate}
\TeacherOnly{%
\begin{tcolorbox}[colback=VE!3,colframe=VE!35,boxrule=.4pt,arc=1mm]
\small\textit{Réponses.}\quad \textbf{1.} Non. \quad \textbf{2.} $5$. \quad \textbf{3.} Entrée; sortie. \quad \textbf{4.} $(25,5)$. \quad \textbf{5.} $64$.
\end{tcolorbox}}

\section{Semaine 4, séance B -- Boucles et organigrammes}
\begin{enumerate}
  \item Une boucle Tant que peut-elle ne jamais s'arrêter ?
  \item Donner l'instruction qui augmente $k$ de $1$.
  \item Que vaut $17\bmod5$ ?
  \item Que vaut le quotient entier de $17$ par $5$ ?
  \item Résoudre $|x+1|<3$.
\end{enumerate}
\TeacherOnly{%
\begin{tcolorbox}[colback=VE!3,colframe=VE!35,boxrule=.4pt,arc=1mm]
\small\textit{Réponses.}\quad \textbf{1.} Oui, si sa condition reste vraie. \quad \textbf{2.} $k\leftarrow k+1$. \quad \textbf{3.} $2$. \quad \textbf{4.} $3$. \quad \textbf{5.} $x\in]-4;2[$.
\end{tcolorbox}}

\section{Semaine 5, séance A -- Ensembles de nombres}
\begin{enumerate}
  \item Calculer $|{-4}-3|$.
  \item Résoudre $|x|=6$.
  \item Écrire $\dfrac13$ avec une précision de $10^{-2}$.
  \item Calculer $(2^3)^2$.
  \item Nier : « il existe un entier pair supérieur à $10$ ».
\end{enumerate}
\TeacherOnly{%
\begin{tcolorbox}[colback=VE!3,colframe=VE!35,boxrule=.4pt,arc=1mm]
\small\textit{Réponses.}\quad \textbf{1.} $7$. \quad \textbf{2.} $x=-6$ ou $x=6$. \quad \textbf{3.} $0{,}33$ au centième. \quad \textbf{4.} $64$. \quad \textbf{5.} Tout entier supérieur à $10$ est impair.
\end{tcolorbox}}

\section{Semaine 5, séance B -- Fractions et puissances}
\begin{enumerate}
  \item Rationaliser $\dfrac1{\sqrt5}$.
  \item Résoudre $|x+1|<3$.
  \item Donner un encadrement décimal de $\sqrt2$ au dixième.
  \item Calculer $\dfrac{5}{12}+\dfrac14$.
  \item Écrire « $2$ est l'unique solution » avec le symbole $\exists!$.
\end{enumerate}
\TeacherOnly{%
\begin{tcolorbox}[colback=VE!3,colframe=VE!35,boxrule=.4pt,arc=1mm]
\small\textit{Réponses.}\quad \textbf{1.} $\dfrac{\sqrt5}{5}$. \quad \textbf{2.} $x\in]-4;2[$. \quad \textbf{3.} $1{,}4<\sqrt2<1{,}5$. \quad \textbf{4.} $\dfrac23$. \quad \textbf{5.} $\exists!x,\ x=2$.
\end{tcolorbox}}

\section{Semaine 6, séance A -- Racines carrées et ordre}
\begin{enumerate}
  \item Simplifier $\dfrac{42}{56}$.
  \item Calculer $2^{-3}$.
  \item Écrire $0{,}00072$ en notation scientifique.
  \item Encadrer $\sqrt{20}$ entre deux entiers consécutifs.
  \item La négation de $x\neq0$ est-elle $x=0$ ?
\end{enumerate}
\TeacherOnly{%
\begin{tcolorbox}[colback=VE!3,colframe=VE!35,boxrule=.4pt,arc=1mm]
\small\textit{Réponses.}\quad \textbf{1.} $\dfrac34$. \quad \textbf{2.} $\dfrac18$. \quad \textbf{3.} $7{,}2\times10^{-4}$. \quad \textbf{4.} $4<\sqrt{20}<5$. \quad \textbf{5.} Oui.
\end{tcolorbox}}

\section{Semaine 6, séance B -- Intervalles et valeur absolue}
\begin{enumerate}
  \item Traduire $|x-3|\leq2$ par un intervalle.
  \item Calculer $\dfrac34-\dfrac16$.
  \item Simplifier $\sqrt{75}$.
  \item Calculer $3^2\times3^{-4}$.
  \item Dire si « $7$ est pair » est une proposition et donner sa valeur de vérité.
\end{enumerate}
\TeacherOnly{%
\begin{tcolorbox}[colback=VE!3,colframe=VE!35,boxrule=.4pt,arc=1mm]
\small\textit{Réponses.}\quad \textbf{1.} $x\in[1;5]$. \quad \textbf{2.} $\dfrac7{12}$. \quad \textbf{3.} $5\sqrt3$. \quad \textbf{4.} $\dfrac19$. \quad \textbf{5.} Proposition fausse.
\end{tcolorbox}}

\section{Semaine 7, séance A -- Équations de droites}
\begin{enumerate}
  \item Donner une droite parallèle à $y=-2x+1$.
  \item Donner la pente d'une droite perpendiculaire à $y=\dfrac12x+3$.
  \item Transformer $y=3x-5$ en forme cartésienne.
  \item Déterminer l'ordonnée à l'origine de $2x+y-7=0$.
  \item Calculer $3^2\times3^{-4}$.
\end{enumerate}
\TeacherOnly{%
\begin{tcolorbox}[colback=VE!3,colframe=VE!35,boxrule=.4pt,arc=1mm]
\small\textit{Réponses.}\quad \textbf{1.} Par exemple $y=-2x+4$. \quad \textbf{2.} $-2$. \quad \textbf{3.} $3x-y-5=0$. \quad \textbf{4.} $7$. \quad \textbf{5.} $\dfrac19$.
\end{tcolorbox}}

\section{Semaine 7, séance B -- Forme paramétrique d'une droite}
\begin{enumerate}
  \item Les droites $y=4x+1$ et $y=4x-6$ sont-elles parallèles ?
  \item Donner une équation de la droite passant par $(2;1)$ et verticale.
  \item Calculer le milieu de $(1;3)$ et $(5;-1)$.
  \item Calculer la distance entre $(0;0)$ et $(3;4)$.
  \item Donner $]-\infty;3]\cup[1;6[$.
\end{enumerate}
\TeacherOnly{%
\begin{tcolorbox}[colback=VE!3,colframe=VE!35,boxrule=.4pt,arc=1mm]
\small\textit{Réponses.}\quad \textbf{1.} Oui. \quad \textbf{2.} $x=2$. \quad \textbf{3.} $(3;1)$. \quad \textbf{4.} $5$. \quad \textbf{5.} $]-\infty;6[$.
\end{tcolorbox}}

\section{Semaine 8, séance A -- Écritures approchées}
\begin{enumerate}
  \item Rationaliser $\dfrac1{\sqrt5}$.
  \item Résoudre $|x+1|<3$.
  \item Donner un encadrement décimal de $\sqrt2$ au dixième.
  \item Calculer $\dfrac{5}{12}+\dfrac14$.
  \item Nier : « $P$ ou $Q$ ».
\end{enumerate}
\TeacherOnly{%
\begin{tcolorbox}[colback=VE!3,colframe=VE!35,boxrule=.4pt,arc=1mm]
\small\textit{Réponses.}\quad \textbf{1.} $\dfrac{\sqrt5}{5}$. \quad \textbf{2.} $x\in]-4;2[$. \quad \textbf{3.} $1{,}4<\sqrt2<1{,}5$. \quad \textbf{4.} $\dfrac23$. \quad \textbf{5.} « non $P$ et non $Q$ ».
\end{tcolorbox}}

\section{Semaine 8, séance B -- Encadrements et propagation}
\begin{enumerate}
  \item Simplifier $\dfrac{42}{56}$.
  \item Calculer $2^{-3}$.
  \item Écrire $0{,}00072$ en notation scientifique.
  \item Encadrer $\sqrt{20}$ entre deux entiers consécutifs.
  \item Donner la réciproque de $P\Rightarrow Q$.
\end{enumerate}
\TeacherOnly{%
\begin{tcolorbox}[colback=VE!3,colframe=VE!35,boxrule=.4pt,arc=1mm]
\small\textit{Réponses.}\quad \textbf{1.} $\dfrac34$. \quad \textbf{2.} $\dfrac18$. \quad \textbf{3.} $7{,}2\times10^{-4}$. \quad \textbf{4.} $4<\sqrt{20}<5$. \quad \textbf{5.} $Q\Rightarrow P$.
\end{tcolorbox}}

\section{Semaine 9, séance A -- Forme canonique}
\begin{enumerate}
  \item Combien de racines réelles si $\Delta=0$ ?
  \item Vérifier si $2$ est racine de $x^3-3x^2-4x+12$.
  \item Factoriser $x^3-4x$.
  \item Donner les valeurs interdites de $\dfrac{x+1}{x^2-4}$.
  \item Calculer $(2^3)^2$.
\end{enumerate}
\TeacherOnly{%
\begin{tcolorbox}[colback=VE!3,colframe=VE!35,boxrule=.4pt,arc=1mm]
\small\textit{Réponses.}\quad \textbf{1.} Une racine double. \quad \textbf{2.} Oui. \quad \textbf{3.} $x(x-2)(x+2)$. \quad \textbf{4.} $-2$ et $2$. \quad \textbf{5.} $64$.
\end{tcolorbox}}

\section{Semaine 9, séance B -- Discriminant}
\begin{enumerate}
  \item Résoudre $\dfrac{x-1}{x+2}=0$.
  \item Donner le signe de $-2(x-1)^2$.
  \item Résoudre $x^2+4x+4>0$.
  \item Développer $(x-4)^2$.
  \item Résoudre $|x+1|<3$.
\end{enumerate}
\TeacherOnly{%
\begin{tcolorbox}[colback=VE!3,colframe=VE!35,boxrule=.4pt,arc=1mm]
\small\textit{Réponses.}\quad \textbf{1.} $x=1$. \quad \textbf{2.} Négatif ou nul. \quad \textbf{3.} $\R\setminus\{-2\}$. \quad \textbf{4.} $x^2-8x+16$. \quad \textbf{5.} $x\in]-4;2[$.
\end{tcolorbox}}

\section{Semaine 10, séance A -- Résolution du second degré}
\begin{enumerate}
  \item Calculer le discriminant de $x^2-5x+6$.
  \item Résoudre $x^2-5x+6=0$.
  \item Factoriser $2x^2+4x+2$.
  \item Donner la forme canonique de $x^2-6x+5$.
  \item Calculer $\dfrac{5}{12}+\dfrac14$.
\end{enumerate}
\TeacherOnly{%
\begin{tcolorbox}[colback=VE!3,colframe=VE!35,boxrule=.4pt,arc=1mm]
\small\textit{Réponses.}\quad \textbf{1.} $\Delta=1$. \quad \textbf{2.} $x=2$ ou $x=3$. \quad \textbf{3.} $2(x+1)^2$. \quad \textbf{4.} $(x-3)^2-4$. \quad \textbf{5.} $\dfrac23$.
\end{tcolorbox}}

\section{Semaine 10, séance B -- Factorisation du trinôme}
\begin{enumerate}
  \item Résoudre $(x-3)(x+4)\leq0$.
  \item Donner le signe de $(x-2)(x+1)$ pour $x>2$.
  \item Résoudre $x^2-9=0$.
  \item Factoriser $x^2+x-6$.
  \item Calculer $2^{-3}$.
\end{enumerate}
\TeacherOnly{%
\begin{tcolorbox}[colback=VE!3,colframe=VE!35,boxrule=.4pt,arc=1mm]
\small\textit{Réponses.}\quad \textbf{1.} $x\in[-4;3]$. \quad \textbf{2.} Positif. \quad \textbf{3.} $x=-3$ ou $x=3$. \quad \textbf{4.} $(x+3)(x-2)$. \quad \textbf{5.} $\dfrac18$.
\end{tcolorbox}}

\section{Semaine 11, séance A -- Signe du trinôme}
\begin{enumerate}
  \item Donner le sommet de $x^2-4x+1$.
  \item Combien de racines réelles si $\Delta<0$ ?
  \item Combien de racines réelles si $\Delta=0$ ?
  \item Vérifier si $2$ est racine de $x^3-3x^2-4x+12$.
  \item Encadrer $\sqrt{20}$ entre deux entiers consécutifs.
\end{enumerate}
\TeacherOnly{%
\begin{tcolorbox}[colback=VE!3,colframe=VE!35,boxrule=.4pt,arc=1mm]
\small\textit{Réponses.}\quad \textbf{1.} $S(2;-3)$. \quad \textbf{2.} Aucune. \quad \textbf{3.} Une racine double. \quad \textbf{4.} Oui. \quad \textbf{5.} $4<\sqrt{20}<5$.
\end{tcolorbox}}

\section{Semaine 11, séance B -- Inéquations du second degré}
\begin{enumerate}
  \item Factoriser $x^3-4x$.
  \item Donner les valeurs interdites de $\dfrac{x+1}{x^2-4}$.
  \item Résoudre $\dfrac{x-1}{x+2}=0$.
  \item Donner le signe de $-2(x-1)^2$.
  \item Calculer $\dfrac34-\dfrac16$.
\end{enumerate}
\TeacherOnly{%
\begin{tcolorbox}[colback=VE!3,colframe=VE!35,boxrule=.4pt,arc=1mm]
\small\textit{Réponses.}\quad \textbf{1.} $x(x-2)(x+2)$. \quad \textbf{2.} $-2$ et $2$. \quad \textbf{3.} $x=1$. \quad \textbf{4.} Négatif ou nul. \quad \textbf{5.} $\dfrac7{12}$.
\end{tcolorbox}}

\section{Semaine 12, séance A -- Polynôme avec une racine connue}
\begin{enumerate}
  \item Résoudre $x^2+4x+4>0$.
  \item Développer $(x-4)^2$.
  \item Calculer le discriminant de $x^2-5x+6$.
  \item Résoudre $x^2-5x+6=0$.
  \item Calculer $3^2\times3^{-4}$.
\end{enumerate}
\TeacherOnly{%
\begin{tcolorbox}[colback=VE!3,colframe=VE!35,boxrule=.4pt,arc=1mm]
\small\textit{Réponses.}\quad \textbf{1.} $\R\setminus\{-2\}$. \quad \textbf{2.} $x^2-8x+16$. \quad \textbf{3.} $\Delta=1$. \quad \textbf{4.} $x=2$ ou $x=3$. \quad \textbf{5.} $\dfrac19$.
\end{tcolorbox}}

\section{Semaine 12, séance B -- Division et identification}
\begin{enumerate}
  \item Factoriser $2x^2+4x+2$.
  \item Donner la forme canonique de $x^2-6x+5$.
  \item Résoudre $(x-3)(x+4)\leq0$.
  \item Donner le signe de $(x-2)(x+1)$ pour $x>2$.
  \item Donner $]-\infty;3]\cup[1;6[$.
\end{enumerate}
\TeacherOnly{%
\begin{tcolorbox}[colback=VE!3,colframe=VE!35,boxrule=.4pt,arc=1mm]
\small\textit{Réponses.}\quad \textbf{1.} $2(x+1)^2$. \quad \textbf{2.} $(x-3)^2-4$. \quad \textbf{3.} $x\in[-4;3]$. \quad \textbf{4.} Positif. \quad \textbf{5.} $]-\infty;6[$.
\end{tcolorbox}}

\section{Semaine 13, séance A -- Signe d'un produit}
\begin{enumerate}
  \item Résoudre $x^2-9=0$.
  \item Factoriser $x^2+x-6$.
  \item Donner le sommet de $x^2-4x+1$.
  \item Combien de racines réelles si $\Delta<0$ ?
  \item Comparer $\sqrt{11}$ et $\dfrac72$.
\end{enumerate}
\TeacherOnly{%
\begin{tcolorbox}[colback=VE!3,colframe=VE!35,boxrule=.4pt,arc=1mm]
\small\textit{Réponses.}\quad \textbf{1.} $x=-3$ ou $x=3$. \quad \textbf{2.} $(x+3)(x-2)$. \quad \textbf{3.} $S(2;-3)$. \quad \textbf{4.} Aucune. \quad \textbf{5.} $\sqrt{11}<\dfrac72$.
\end{tcolorbox}}

\section{Semaine 13, séance B -- Équations rationnelles}
\begin{enumerate}
  \item Combien de racines réelles si $\Delta=0$ ?
  \item Vérifier si $2$ est racine de $x^3-3x^2-4x+12$.
  \item Factoriser $x^3-4x$.
  \item Donner les valeurs interdites de $\dfrac{x+1}{x^2-4}$.
  \item Résoudre $|x|=6$.
\end{enumerate}
\TeacherOnly{%
\begin{tcolorbox}[colback=VE!3,colframe=VE!35,boxrule=.4pt,arc=1mm]
\small\textit{Réponses.}\quad \textbf{1.} Une racine double. \quad \textbf{2.} Oui. \quad \textbf{3.} $x(x-2)(x+2)$. \quad \textbf{4.} $-2$ et $2$. \quad \textbf{5.} $x=-6$ ou $x=6$.
\end{tcolorbox}}

\section{Semaine 14, séance A -- Équations de cercles}
\begin{enumerate}
  \item Si $d(\Omega,\Delta)<r$, quelle est la position de la droite ?
  \item La tangente en $A$ est perpendiculaire à quel segment ?
  \item Calculer $d((0;0),(6;8))$.
  \item Donner l'équation du cercle de diamètre $[AB]$, avec $A(-2;0)$ et $B(2;0)$.
  \item Calculer $(2^3)^2$.
\end{enumerate}
\TeacherOnly{%
\begin{tcolorbox}[colback=VE!3,colframe=VE!35,boxrule=.4pt,arc=1mm]
\small\textit{Réponses.}\quad \textbf{1.} Sécante. \quad \textbf{2.} Au rayon $[\Omega A]$. \quad \textbf{3.} $10$. \quad \textbf{4.} $x^2+y^2=4$. \quad \textbf{5.} $64$.
\end{tcolorbox}}

\section{Semaine 14, séance B -- Intersection droite-cercle}
\begin{enumerate}
  \item Deux cercles ont $r_1=3$, $r_2=2$, $d=5$. Leur position ?
  \item Deux cercles ont $r_1=5$, $r_2=2$, $d=1$. Leur position ?
  \item Compléter le carré : $x^2-6x=(x-3)^2-\ldots$
  \item Donner une équation de la tangente à $x^2+y^2=25$ en $(5;0)$.
  \item Résoudre $|x+1|<3$.
\end{enumerate}
\TeacherOnly{%
\begin{tcolorbox}[colback=VE!3,colframe=VE!35,boxrule=.4pt,arc=1mm]
\small\textit{Réponses.}\quad \textbf{1.} Tangents extérieurement. \quad \textbf{2.} L'un est intérieur à l'autre. \quad \textbf{3.} $9$. \quad \textbf{4.} $x=5$. \quad \textbf{5.} $x\in]-4;2[$.
\end{tcolorbox}}

\section{Semaine 15, séance A -- Domaine et image d'une fonction}
\begin{enumerate}
  \item Résoudre $|x|=5$.
  \item La fonction cube est-elle croissante sur $\R$ ?
  \item Calculer l'image de $-2$ par $x\mapsto x^3$.
  \item Donner le domaine de $x\mapsto\dfrac1{x-4}$.
  \item Calculer $\dfrac{5}{12}+\dfrac14$.
\end{enumerate}
\TeacherOnly{%
\begin{tcolorbox}[colback=VE!3,colframe=VE!35,boxrule=.4pt,arc=1mm]
\small\textit{Réponses.}\quad \textbf{1.} $x=-5$ ou $x=5$. \quad \textbf{2.} Oui. \quad \textbf{3.} $-8$. \quad \textbf{4.} $\R\setminus\{4\}$. \quad \textbf{5.} $\dfrac23$.
\end{tcolorbox}}

\section{Semaine 15, séance B -- Lecture graphique}
\begin{enumerate}
  \item Résoudre $\sqrt{x}=3$.
  \item Comparer $(-3)^2$ et $(-2)^2$.
  \item Si $a<b$ et $f$ est croissante, comparer $f(a)$ et $f(b)$.
  \item Donner le sommet de $x\mapsto-(x+1)^2+4$.
  \item Calculer $2^{-3}$.
\end{enumerate}
\TeacherOnly{%
\begin{tcolorbox}[colback=VE!3,colframe=VE!35,boxrule=.4pt,arc=1mm]
\small\textit{Réponses.}\quad \textbf{1.} $x=9$. \quad \textbf{2.} $9>4$. \quad \textbf{3.} $f(a)\leq f(b)$. \quad \textbf{4.} $(-1;4)$. \quad \textbf{5.} $\dfrac18$.
\end{tcolorbox}}

\section{Semaine 16, séance A -- Variations et extremum}
\begin{enumerate}
  \item Résoudre $|x-2|\leq3$.
  \item Calculer $f(0)$ pour $f(x)=\sqrt{x+4}$.
  \item La fonction racine carrée est-elle définie en $-1$ ?
  \item Donner les zéros de $(x-1)^2-9$.
  \item Encadrer $\sqrt{20}$ entre deux entiers consécutifs.
\end{enumerate}
\TeacherOnly{%
\begin{tcolorbox}[colback=VE!3,colframe=VE!35,boxrule=.4pt,arc=1mm]
\small\textit{Réponses.}\quad \textbf{1.} $x\in[-1;5]$. \quad \textbf{2.} $2$. \quad \textbf{3.} Non. \quad \textbf{4.} $-2$ et $4$. \quad \textbf{5.} $4<\sqrt{20}<5$.
\end{tcolorbox}}

\section{Semaine 16, séance B -- Parité}
\begin{enumerate}
  \item Une fonction impaire vérifie quelle relation ?
  \item Donner l'image de $3$ par $x\mapsto|x-5|$.
  \item Si $f(x)=x^2-2x$, calculer $f(-1)$.
  \item Donner le domaine de $x\mapsto\sqrt{x-3}$.
  \item Calculer $\dfrac34-\dfrac16$.
\end{enumerate}
\TeacherOnly{%
\begin{tcolorbox}[colback=VE!3,colframe=VE!35,boxrule=.4pt,arc=1mm]
\small\textit{Réponses.}\quad \textbf{1.} $f(-x)=-f(x)$. \quad \textbf{2.} $2$. \quad \textbf{3.} $3$. \quad \textbf{4.} $[3;+\infty[$. \quad \textbf{5.} $\dfrac7{12}$.
\end{tcolorbox}}

\section{Semaine 17, séance A -- Fonction carré}
\begin{enumerate}
  \item La fonction carré est-elle paire ?
  \item Donner le minimum de $(x-2)^2+1$.
  \item Résoudre $|x|=5$.
  \item La fonction cube est-elle croissante sur $\R$ ?
  \item Calculer $3^2\times3^{-4}$.
\end{enumerate}
\TeacherOnly{%
\begin{tcolorbox}[colback=VE!3,colframe=VE!35,boxrule=.4pt,arc=1mm]
\small\textit{Réponses.}\quad \textbf{1.} Oui. \quad \textbf{2.} $1$, atteint en $2$. \quad \textbf{3.} $x=-5$ ou $x=5$. \quad \textbf{4.} Oui. \quad \textbf{5.} $\dfrac19$.
\end{tcolorbox}}

\section{Semaine 17, séance B -- Fonction cube}
\begin{enumerate}
  \item Calculer l'image de $-2$ par $x\mapsto x^3$.
  \item Donner le domaine de $x\mapsto\dfrac1{x-4}$.
  \item Résoudre $\sqrt{x}=3$.
  \item Comparer $(-3)^2$ et $(-2)^2$.
  \item Donner $]-\infty;3]\cup[1;6[$.
\end{enumerate}
\TeacherOnly{%
\begin{tcolorbox}[colback=VE!3,colframe=VE!35,boxrule=.4pt,arc=1mm]
\small\textit{Réponses.}\quad \textbf{1.} $-8$. \quad \textbf{2.} $\R\setminus\{4\}$. \quad \textbf{3.} $x=9$. \quad \textbf{4.} $9>4$. \quad \textbf{5.} $]-\infty;6[$.
\end{tcolorbox}}

\section{Semaine 18, séance A -- Valeur absolue}
\begin{enumerate}
  \item Si $a<b$ et $f$ est croissante, comparer $f(a)$ et $f(b)$.
  \item Donner le sommet de $x\mapsto-(x+1)^2+4$.
  \item Résoudre $|x-2|\leq3$.
  \item Calculer $f(0)$ pour $f(x)=\sqrt{x+4}$.
  \item Comparer $\sqrt{11}$ et $\dfrac72$.
\end{enumerate}
\TeacherOnly{%
\begin{tcolorbox}[colback=VE!3,colframe=VE!35,boxrule=.4pt,arc=1mm]
\small\textit{Réponses.}\quad \textbf{1.} $f(a)\leq f(b)$. \quad \textbf{2.} $(-1;4)$. \quad \textbf{3.} $x\in[-1;5]$. \quad \textbf{4.} $2$. \quad \textbf{5.} $\sqrt{11}<\dfrac72$.
\end{tcolorbox}}

\section{Semaine 18, séance B -- Racine carrée}
\begin{enumerate}
  \item La fonction racine carrée est-elle définie en $-1$ ?
  \item Donner les zéros de $(x-1)^2-9$.
  \item Une fonction impaire vérifie quelle relation ?
  \item Donner l'image de $3$ par $x\mapsto|x-5|$.
  \item Résoudre $|x|=6$.
\end{enumerate}
\TeacherOnly{%
\begin{tcolorbox}[colback=VE!3,colframe=VE!35,boxrule=.4pt,arc=1mm]
\small\textit{Réponses.}\quad \textbf{1.} Non. \quad \textbf{2.} $-2$ et $4$. \quad \textbf{3.} $f(-x)=-f(x)$. \quad \textbf{4.} $2$. \quad \textbf{5.} $x=-6$ ou $x=6$.
\end{tcolorbox}}

\section{Semaine 19, séance A -- Degrés et radians}
\begin{enumerate}
  \item Convertir $180^\circ$ en radians.
  \item Convertir $\dfrac{\pi}{3}$ en degrés.
  \item Donner $\cos0$ et $\sin0$.
  \item Donner $\cos\dfrac{\pi}{2}$ et $\sin\dfrac{\pi}{2}$.
  \item Calculer $(2^3)^2$.
\end{enumerate}
\TeacherOnly{%
\begin{tcolorbox}[colback=VE!3,colframe=VE!35,boxrule=.4pt,arc=1mm]
\small\textit{Réponses.}\quad \textbf{1.} $\pi$. \quad \textbf{2.} $60^\circ$. \quad \textbf{3.} $1$ et $0$. \quad \textbf{4.} $0$ et $1$. \quad \textbf{5.} $64$.
\end{tcolorbox}}

\section{Semaine 19, séance B -- Cercle trigonométrique}
\begin{enumerate}
  \item Compléter : $\cos^2x+\sin^2x=\ldots$
  \item Donner $\cos\dfrac{\pi}{3}$.
  \item Donner $\sin\dfrac{\pi}{6}$.
  \item Réduire $\dfrac{13\pi}{6}$ modulo $2\pi$.
  \item Résoudre $|x+1|<3$.
\end{enumerate}
\TeacherOnly{%
\begin{tcolorbox}[colback=VE!3,colframe=VE!35,boxrule=.4pt,arc=1mm]
\small\textit{Réponses.}\quad \textbf{1.} $1$. \quad \textbf{2.} $\dfrac12$. \quad \textbf{3.} $\dfrac12$. \quad \textbf{4.} $\dfrac{\pi}{6}$. \quad \textbf{5.} $x\in]-4;2[$.
\end{tcolorbox}}

\section{Semaine 20, séance A -- Angles remarquables}
\begin{enumerate}
  \item Quelle est la parité du cosinus ?
  \item Quelle est la parité du sinus ?
  \item Convertir $150^\circ$ en radians.
  \item Donner $\cos\pi$ et $\sin\pi$.
  \item Calculer $\dfrac{5}{12}+\dfrac14$.
\end{enumerate}
\TeacherOnly{%
\begin{tcolorbox}[colback=VE!3,colframe=VE!35,boxrule=.4pt,arc=1mm]
\small\textit{Réponses.}\quad \textbf{1.} Paire. \quad \textbf{2.} Impaire. \quad \textbf{3.} $\dfrac{5\pi}{6}$. \quad \textbf{4.} $-1$ et $0$. \quad \textbf{5.} $\dfrac23$.
\end{tcolorbox}}

\section{Semaine 20, séance B -- Relation fondamentale}
\begin{enumerate}
  \item Si $\sin x=\dfrac35$ et $x$ est dans le premier quadrant, donner $\cos x$.
  \item Donner une mesure positive de l'angle $-\dfrac{\pi}{2}$ modulo $2\pi$.
  \item Dans quel quadrant se trouve $\dfrac{3\pi}{4}$ ?
  \item Donner la période du sinus.
  \item Calculer $2^{-3}$.
\end{enumerate}
\TeacherOnly{%
\begin{tcolorbox}[colback=VE!3,colframe=VE!35,boxrule=.4pt,arc=1mm]
\small\textit{Réponses.}\quad \textbf{1.} $\dfrac45$. \quad \textbf{2.} $\dfrac{3\pi}{2}$. \quad \textbf{3.} Deuxième quadrant. \quad \textbf{4.} $2\pi$. \quad \textbf{5.} $\dfrac18$.
\end{tcolorbox}}

\section{Semaine 21, séance A -- Parité et périodicité}
\begin{enumerate}
  \item Convertir $180^\circ$ en radians.
  \item Convertir $\dfrac{\pi}{3}$ en degrés.
  \item Donner $\cos0$ et $\sin0$.
  \item Donner $\cos\dfrac{\pi}{2}$ et $\sin\dfrac{\pi}{2}$.
  \item Encadrer $\sqrt{20}$ entre deux entiers consécutifs.
\end{enumerate}
\TeacherOnly{%
\begin{tcolorbox}[colback=VE!3,colframe=VE!35,boxrule=.4pt,arc=1mm]
\small\textit{Réponses.}\quad \textbf{1.} $\pi$. \quad \textbf{2.} $60^\circ$. \quad \textbf{3.} $1$ et $0$. \quad \textbf{4.} $0$ et $1$. \quad \textbf{5.} $4<\sqrt{20}<5$.
\end{tcolorbox}}

\section{Semaine 21, séance B -- Fonctions et trigonométrie}
\begin{enumerate}
  \item Compléter : $\cos^2x+\sin^2x=\ldots$
  \item Donner $\cos\dfrac{\pi}{3}$.
  \item Donner $\sin\dfrac{\pi}{6}$.
  \item Réduire $\dfrac{13\pi}{6}$ modulo $2\pi$.
  \item Calculer $\dfrac34-\dfrac16$.
\end{enumerate}
\TeacherOnly{%
\begin{tcolorbox}[colback=VE!3,colframe=VE!35,boxrule=.4pt,arc=1mm]
\small\textit{Réponses.}\quad \textbf{1.} $1$. \quad \textbf{2.} $\dfrac12$. \quad \textbf{3.} $\dfrac12$. \quad \textbf{4.} $\dfrac{\pi}{6}$. \quad \textbf{5.} $\dfrac7{12}$.
\end{tcolorbox}}

\section{Semaine 22, séance A -- Courbes de référence}
\begin{enumerate}
  \item Calculer l'image de $-2$ par $x\mapsto x^3$.
  \item Donner le domaine de $x\mapsto\dfrac1{x-4}$.
  \item Résoudre $\sqrt{x}=3$.
  \item Comparer $(-3)^2$ et $(-2)^2$.
  \item Calculer $3^2\times3^{-4}$.
\end{enumerate}
\TeacherOnly{%
\begin{tcolorbox}[colback=VE!3,colframe=VE!35,boxrule=.4pt,arc=1mm]
\small\textit{Réponses.}\quad \textbf{1.} $-8$. \quad \textbf{2.} $\R\setminus\{4\}$. \quad \textbf{3.} $x=9$. \quad \textbf{4.} $9>4$. \quad \textbf{5.} $\dfrac19$.
\end{tcolorbox}}

\section{Semaine 22, séance B -- Tableaux de variation}
\begin{enumerate}
  \item Si $a<b$ et $f$ est croissante, comparer $f(a)$ et $f(b)$.
  \item Donner le sommet de $x\mapsto-(x+1)^2+4$.
  \item Résoudre $|x-2|\leq3$.
  \item Calculer $f(0)$ pour $f(x)=\sqrt{x+4}$.
  \item Donner $]-\infty;3]\cup[1;6[$.
\end{enumerate}
\TeacherOnly{%
\begin{tcolorbox}[colback=VE!3,colframe=VE!35,boxrule=.4pt,arc=1mm]
\small\textit{Réponses.}\quad \textbf{1.} $f(a)\leq f(b)$. \quad \textbf{2.} $(-1;4)$. \quad \textbf{3.} $x\in[-1;5]$. \quad \textbf{4.} $2$. \quad \textbf{5.} $]-\infty;6[$.
\end{tcolorbox}}

\section{Semaine 23, séance A -- Problèmes sur les fonctions}
\begin{enumerate}
  \item La fonction racine carrée est-elle définie en $-1$ ?
  \item Donner les zéros de $(x-1)^2-9$.
  \item Une fonction impaire vérifie quelle relation ?
  \item Donner l'image de $3$ par $x\mapsto|x-5|$.
  \item Comparer $\sqrt{11}$ et $\dfrac72$.
\end{enumerate}
\TeacherOnly{%
\begin{tcolorbox}[colback=VE!3,colframe=VE!35,boxrule=.4pt,arc=1mm]
\small\textit{Réponses.}\quad \textbf{1.} Non. \quad \textbf{2.} $-2$ et $4$. \quad \textbf{3.} $f(-x)=-f(x)$. \quad \textbf{4.} $2$. \quad \textbf{5.} $\sqrt{11}<\dfrac72$.
\end{tcolorbox}}

\section{Semaine 23, séance B -- Pratique mixte}
\begin{enumerate}
  \item Calculer $(2^3)^2$.
  \item Combien de racines réelles si $\Delta=0$ ?
  \item Donner le sommet de $x\mapsto-(x+1)^2+4$.
  \item Calculer $(1;3)\cdot(-3;1)$.
  \item La somme des fréquences d'une série vaut combien ?
\end{enumerate}
\TeacherOnly{%
\begin{tcolorbox}[colback=VE!3,colframe=VE!35,boxrule=.4pt,arc=1mm]
\small\textit{Réponses.}\quad \textbf{1.} $64$. \quad \textbf{2.} Une racine double. \quad \textbf{3.} $(-1;4)$. \quad \textbf{4.} $0$. \quad \textbf{5.} $1$ ou $100\%$.
\end{tcolorbox}}

\section{Semaine 24, séance A -- Vecteurs : direction et norme}
\begin{enumerate}
  \item Les vecteurs $(1;2)$ et $(-4;2)$ sont-ils orthogonaux ?
  \item Calculer la norme de $(3;4)$.
  \item Compléter : $\vv{AB}+\vv{BC}=\ldots$
  \item Donner un vecteur colinéaire à $(3;-2)$.
  \item Calculer $(2^3)^2$.
\end{enumerate}
\TeacherOnly{%
\begin{tcolorbox}[colback=VE!3,colframe=VE!35,boxrule=.4pt,arc=1mm]
\small\textit{Réponses.}\quad \textbf{1.} Oui. \quad \textbf{2.} $5$. \quad \textbf{3.} $\vv{AC}$. \quad \textbf{4.} Par exemple $(6;-4)$. \quad \textbf{5.} $64$.
\end{tcolorbox}}

\section{Semaine 24, séance B -- Égalité et opérations}
\begin{enumerate}
  \item Calculer le déterminant de $(2;1)$ et $(3;4)$.
  \item Si $\vec u\cdot\vec v=0$ et les vecteurs sont non nuls, que conclure ?
  \item Si $\vv{AB}=\vv{DC}$, quelle figure est $ABCD$ ?
  \item Donner les coordonnées du centre de gravité de $(0;0)$, $(3;0)$, $(0;6)$.
  \item Résoudre $|x+1|<3$.
\end{enumerate}
\TeacherOnly{%
\begin{tcolorbox}[colback=VE!3,colframe=VE!35,boxrule=.4pt,arc=1mm]
\small\textit{Réponses.}\quad \textbf{1.} $5$. \quad \textbf{2.} Ils sont orthogonaux. \quad \textbf{3.} Un parallélogramme. \quad \textbf{4.} $(1;2)$. \quad \textbf{5.} $x\in]-4;2[$.
\end{tcolorbox}}

\section{Semaine 25, séance A -- Coordonnées et milieu}
\begin{enumerate}
  \item Calculer $(1;3)\cdot(-3;1)$.
  \item Les points $(0;0)$, $(2;3)$, $(4;6)$ sont-ils alignés ?
  \item Donner $\vv{BA}$ si $\vv{AB}=(5;-2)$.
  \item Calculer $\|(1;-2)\|^2$.
  \item Calculer $\dfrac{5}{12}+\dfrac14$.
\end{enumerate}
\TeacherOnly{%
\begin{tcolorbox}[colback=VE!3,colframe=VE!35,boxrule=.4pt,arc=1mm]
\small\textit{Réponses.}\quad \textbf{1.} $0$. \quad \textbf{2.} Oui. \quad \textbf{3.} $(-5;2)$. \quad \textbf{4.} $5$. \quad \textbf{5.} $\dfrac23$.
\end{tcolorbox}}

\section{Semaine 25, séance B -- Colinéarité}
\begin{enumerate}
  \item Que vaut $\vec u\cdot\vec u$ ?
  \item Donner le point $D$ tel que $ABCD$ soit un parallélogramme, avec $A(0;0)$, $B(2;0)$, $C(3;1)$.
  \item Pour $A(1;2)$ et $B(4;-2)$, donner $\vv{AB}$.
  \item Donner le milieu de $A(1;2)$ et $B(4;-2)$.
  \item Calculer $2^{-3}$.
\end{enumerate}
\TeacherOnly{%
\begin{tcolorbox}[colback=VE!3,colframe=VE!35,boxrule=.4pt,arc=1mm]
\small\textit{Réponses.}\quad \textbf{1.} $\|\vec u\|^2$. \quad \textbf{2.} $D(1;1)$. \quad \textbf{3.} $(3;-4)$. \quad \textbf{4.} $\left(\dfrac52;0\right)$. \quad \textbf{5.} $\dfrac18$.
\end{tcolorbox}}

\section{Semaine 26, séance A -- Alignement et parallélisme}
\begin{enumerate}
  \item Les vecteurs $(2;3)$ et $(4;6)$ sont-ils colinéaires ?
  \item Calculer $(2;-1)\cdot(3;4)$.
  \item Les vecteurs $(1;2)$ et $(-4;2)$ sont-ils orthogonaux ?
  \item Calculer la norme de $(3;4)$.
  \item Encadrer $\sqrt{20}$ entre deux entiers consécutifs.
\end{enumerate}
\TeacherOnly{%
\begin{tcolorbox}[colback=VE!3,colframe=VE!35,boxrule=.4pt,arc=1mm]
\small\textit{Réponses.}\quad \textbf{1.} Oui. \quad \textbf{2.} $2$. \quad \textbf{3.} Oui. \quad \textbf{4.} $5$. \quad \textbf{5.} $4<\sqrt{20}<5$.
\end{tcolorbox}}

\section{Semaine 26, séance B -- Produit scalaire}
\begin{enumerate}
  \item Compléter : $\vv{AB}+\vv{BC}=\ldots$
  \item Donner un vecteur colinéaire à $(3;-2)$.
  \item Calculer le déterminant de $(2;1)$ et $(3;4)$.
  \item Si $\vec u\cdot\vec v=0$ et les vecteurs sont non nuls, que conclure ?
  \item Calculer $\dfrac34-\dfrac16$.
\end{enumerate}
\TeacherOnly{%
\begin{tcolorbox}[colback=VE!3,colframe=VE!35,boxrule=.4pt,arc=1mm]
\small\textit{Réponses.}\quad \textbf{1.} $\vv{AC}$. \quad \textbf{2.} Par exemple $(6;-4)$. \quad \textbf{3.} $5$. \quad \textbf{4.} Ils sont orthogonaux. \quad \textbf{5.} $\dfrac7{12}$.
\end{tcolorbox}}

\section{Semaine 27, séance A -- Orthogonalité}
\begin{enumerate}
  \item Si $\vv{AB}=\vv{DC}$, quelle figure est $ABCD$ ?
  \item Donner les coordonnées du centre de gravité de $(0;0)$, $(3;0)$, $(0;6)$.
  \item Calculer $(1;3)\cdot(-3;1)$.
  \item Les points $(0;0)$, $(2;3)$, $(4;6)$ sont-ils alignés ?
  \item Calculer $3^2\times3^{-4}$.
\end{enumerate}
\TeacherOnly{%
\begin{tcolorbox}[colback=VE!3,colframe=VE!35,boxrule=.4pt,arc=1mm]
\small\textit{Réponses.}\quad \textbf{1.} Un parallélogramme. \quad \textbf{2.} $(1;2)$. \quad \textbf{3.} $0$. \quad \textbf{4.} Oui. \quad \textbf{5.} $\dfrac19$.
\end{tcolorbox}}

\section{Semaine 27, séance B -- Problèmes vectoriels}
\begin{enumerate}
  \item Donner $\vv{BA}$ si $\vv{AB}=(5;-2)$.
  \item Calculer $\|(1;-2)\|^2$.
  \item Que vaut $\vec u\cdot\vec u$ ?
  \item Donner le point $D$ tel que $ABCD$ soit un parallélogramme, avec $A(0;0)$, $B(2;0)$, $C(3;1)$.
  \item Donner $]-\infty;3]\cup[1;6[$.
\end{enumerate}
\TeacherOnly{%
\begin{tcolorbox}[colback=VE!3,colframe=VE!35,boxrule=.4pt,arc=1mm]
\small\textit{Réponses.}\quad \textbf{1.} $(-5;2)$. \quad \textbf{2.} $5$. \quad \textbf{3.} $\|\vec u\|^2$. \quad \textbf{4.} $D(1;1)$. \quad \textbf{5.} $]-\infty;6[$.
\end{tcolorbox}}

\section{Semaine 28, séance A -- Population et variable}
\begin{enumerate}
  \item Calculer la moyenne de $2,4,6,8$.
  \item Donner la médiane de $1,3,3,7,9$.
  \item Donner la médiane de $1,3,7,9$.
  \item Calculer la fréquence de $12$ élèves parmi $30$.
  \item Comparer $\sqrt{11}$ et $\dfrac72$.
\end{enumerate}
\TeacherOnly{%
\begin{tcolorbox}[colback=VE!3,colframe=VE!35,boxrule=.4pt,arc=1mm]
\small\textit{Réponses.}\quad \textbf{1.} $5$. \quad \textbf{2.} $3$. \quad \textbf{3.} $5$. \quad \textbf{4.} $0{,}4=40\%$. \quad \textbf{5.} $\sqrt{11}<\dfrac72$.
\end{tcolorbox}}

\section{Semaine 28, séance B -- Effectifs et fréquences}
\begin{enumerate}
  \item Quelle est la variance de $5,5,5,5$ ?
  \item L'écart-type peut-il être négatif ?
  \item Quel graphique convient à des données continues regroupées en classes ?
  \item Si on ajoute $3$ à toutes les valeurs, comment change la moyenne ?
  \item Résoudre $|x|=6$.
\end{enumerate}
\TeacherOnly{%
\begin{tcolorbox}[colback=VE!3,colframe=VE!35,boxrule=.4pt,arc=1mm]
\small\textit{Réponses.}\quad \textbf{1.} $0$. \quad \textbf{2.} Non. \quad \textbf{3.} Un histogramme. \quad \textbf{4.} Elle augmente de $3$. \quad \textbf{5.} $x=-6$ ou $x=6$.
\end{tcolorbox}}

\section{Semaine 29, séance A -- Moyenne et médiane}
\begin{enumerate}
  \item Donner l'étendue de $4,7,9,12$.
  \item Donner le mode de $2,3,3,3,5,6$.
  \item Calculer la moyenne pondérée de $10$ (coef. $2$) et $16$ (coef. $1$).
  \item Si toutes les valeurs sont multipliées par $2$, comment change l'étendue ?
  \item Calculer $(2^3)^2$.
\end{enumerate}
\TeacherOnly{%
\begin{tcolorbox}[colback=VE!3,colframe=VE!35,boxrule=.4pt,arc=1mm]
\small\textit{Réponses.}\quad \textbf{1.} $8$. \quad \textbf{2.} $3$. \quad \textbf{3.} $12$. \quad \textbf{4.} Elle est multipliée par $2$. \quad \textbf{5.} $64$.
\end{tcolorbox}}

\section{Semaine 29, séance B -- Variance et écart-type}
\begin{enumerate}
  \item Quel quartile correspond à la médiane ?
  \item La somme des fréquences d'une série vaut combien ?
  \item Donner l'effectif total des effectifs $4,7,5,4$.
  \item Une fréquence de $0{,}35$ correspond à quel pourcentage ?
  \item Résoudre $|x+1|<3$.
\end{enumerate}
\TeacherOnly{%
\begin{tcolorbox}[colback=VE!3,colframe=VE!35,boxrule=.4pt,arc=1mm]
\small\textit{Réponses.}\quad \textbf{1.} $Q_2$. \quad \textbf{2.} $1$ ou $100\%$. \quad \textbf{3.} $20$. \quad \textbf{4.} $35\%$. \quad \textbf{5.} $x\in]-4;2[$.
\end{tcolorbox}}

\section{Semaine 30, séance A -- Diagrammes en bâtons et circulaires}
\begin{enumerate}
  \item Quel indicateur est le plus sensible à une valeur extrême : moyenne ou médiane ?
  \item Une série a variance $9$. Donner son écart-type.
  \item Calculer la moyenne de $2,4,6,8$.
  \item Donner la médiane de $1,3,3,7,9$.
  \item Calculer $\dfrac{5}{12}+\dfrac14$.
\end{enumerate}
\TeacherOnly{%
\begin{tcolorbox}[colback=VE!3,colframe=VE!35,boxrule=.4pt,arc=1mm]
\small\textit{Réponses.}\quad \textbf{1.} La moyenne. \quad \textbf{2.} $3$. \quad \textbf{3.} $5$. \quad \textbf{4.} $3$. \quad \textbf{5.} $\dfrac23$.
\end{tcolorbox}}

\section{Semaine 30, séance B -- Histogramme}
\begin{enumerate}
  \item Donner la médiane de $1,3,7,9$.
  \item Calculer la fréquence de $12$ élèves parmi $30$.
  \item Quelle est la variance de $5,5,5,5$ ?
  \item L'écart-type peut-il être négatif ?
  \item Calculer $2^{-3}$.
\end{enumerate}
\TeacherOnly{%
\begin{tcolorbox}[colback=VE!3,colframe=VE!35,boxrule=.4pt,arc=1mm]
\small\textit{Réponses.}\quad \textbf{1.} $5$. \quad \textbf{2.} $0{,}4=40\%$. \quad \textbf{3.} $0$. \quad \textbf{4.} Non. \quad \textbf{5.} $\dfrac18$.
\end{tcolorbox}}

\section{Semaine 31, séance A -- Interpréter une représentation}
\begin{enumerate}
  \item Quel graphique convient à des données continues regroupées en classes ?
  \item Si on ajoute $3$ à toutes les valeurs, comment change la moyenne ?
  \item Donner l'étendue de $4,7,9,12$.
  \item Donner le mode de $2,3,3,3,5,6$.
  \item Encadrer $\sqrt{20}$ entre deux entiers consécutifs.
\end{enumerate}
\TeacherOnly{%
\begin{tcolorbox}[colback=VE!3,colframe=VE!35,boxrule=.4pt,arc=1mm]
\small\textit{Réponses.}\quad \textbf{1.} Un histogramme. \quad \textbf{2.} Elle augmente de $3$. \quad \textbf{3.} $8$. \quad \textbf{4.} $3$. \quad \textbf{5.} $4<\sqrt{20}<5$.
\end{tcolorbox}}

\section{Semaine 31, séance B -- Problème statistique}
\begin{enumerate}
  \item Calculer la moyenne pondérée de $10$ (coef. $2$) et $16$ (coef. $1$).
  \item Si toutes les valeurs sont multipliées par $2$, comment change l'étendue ?
  \item Quel quartile correspond à la médiane ?
  \item La somme des fréquences d'une série vaut combien ?
  \item Calculer $\dfrac34-\dfrac16$.
\end{enumerate}
\TeacherOnly{%
\begin{tcolorbox}[colback=VE!3,colframe=VE!35,boxrule=.4pt,arc=1mm]
\small\textit{Réponses.}\quad \textbf{1.} $12$. \quad \textbf{2.} Elle est multipliée par $2$. \quad \textbf{3.} $Q_2$. \quad \textbf{4.} $1$ ou $100\%$. \quad \textbf{5.} $\dfrac7{12}$.
\end{tcolorbox}}

\section{Semaine 32, séance A -- Révisions mixtes I}
\begin{enumerate}
  \item Simplifier $\sqrt{75}$.
  \item Factoriser $x^2+x-6$.
  \item Résoudre $\sqrt{x}=3$.
  \item Si $\vec u\cdot\vec v=0$ et les vecteurs sont non nuls, que conclure ?
  \item Calculer la moyenne pondérée de $10$ (coef. $2$) et $16$ (coef. $1$).
\end{enumerate}
\TeacherOnly{%
\begin{tcolorbox}[colback=VE!3,colframe=VE!35,boxrule=.4pt,arc=1mm]
\small\textit{Réponses.}\quad \textbf{1.} $5\sqrt3$. \quad \textbf{2.} $(x+3)(x-2)$. \quad \textbf{3.} $x=9$. \quad \textbf{4.} Ils sont orthogonaux. \quad \textbf{5.} $12$.
\end{tcolorbox}}

\section{Semaine 32, séance B -- Révisions mixtes II}
\begin{enumerate}
  \item Donner $]-\infty;3]\cup[1;6[$.
  \item Combien de racines réelles si $\Delta=0$ ?
  \item Donner le sommet de $x\mapsto-(x+1)^2+4$.
  \item Calculer $(1;3)\cdot(-3;1)$.
  \item La somme des fréquences d'une série vaut combien ?
\end{enumerate}
\TeacherOnly{%
\begin{tcolorbox}[colback=VE!3,colframe=VE!35,boxrule=.4pt,arc=1mm]
\small\textit{Réponses.}\quad \textbf{1.} $]-\infty;6[$. \quad \textbf{2.} Une racine double. \quad \textbf{3.} $(-1;4)$. \quad \textbf{4.} $0$. \quad \textbf{5.} $1$ ou $100\%$.
\end{tcolorbox}}

\section{Semaine 33, séance A -- Remédiation ciblée}
\begin{enumerate}
  \item Calculer $|{-4}-3|$.
  \item Donner les valeurs interdites de $\dfrac{x+1}{x^2-4}$.
  \item La fonction racine carrée est-elle définie en $-1$ ?
  \item Calculer $\|(1;-2)\|^2$.
  \item Quel indicateur est le plus sensible à une valeur extrême : moyenne ou médiane ?
\end{enumerate}
\TeacherOnly{%
\begin{tcolorbox}[colback=VE!3,colframe=VE!35,boxrule=.4pt,arc=1mm]
\small\textit{Réponses.}\quad \textbf{1.} $7$. \quad \textbf{2.} $-2$ et $2$. \quad \textbf{3.} Non. \quad \textbf{4.} $5$. \quad \textbf{5.} La moyenne.
\end{tcolorbox}}

\section{Semaine 33, séance B -- Approfondissement et recherche}
\begin{enumerate}
  \item Calculer $(2^3)^2$.
  \item Résoudre $x^2+4x+4>0$.
  \item Donner l'image de $3$ par $x\mapsto|x-5|$.
  \item Pour $A(1;2)$ et $B(4;-2)$, donner $\vv{AB}$.
  \item Donner la médiane de $1,3,3,7,9$.
\end{enumerate}
\TeacherOnly{%
\begin{tcolorbox}[colback=VE!3,colframe=VE!35,boxrule=.4pt,arc=1mm]
\small\textit{Réponses.}\quad \textbf{1.} $64$. \quad \textbf{2.} $\R\setminus\{-2\}$. \quad \textbf{3.} $2$. \quad \textbf{4.} $(3;-4)$. \quad \textbf{5.} $3$.
\end{tcolorbox}}

